\section{Introduction}\label{sec:intro}

A reasoning model that writes out its chain of thought commits to one token at every step \cite{wei2022chain}.
The other candidates at that step are discarded.
Latent reasoning methods relax this commitment by feeding the model continuous states instead of tokens
\cite{hao2024coconut,zhang2025soft}.
A common motivation is \emph{superposition}: a continuous state might hold several candidate paths at once and let
the model explore them in parallel.
In theory this can help, since a two-layer transformer with continuous thoughts can solve graph reachability by
keeping several search frontiers in one state \cite{zhu2025superposition}.
\ST{} \cite{zhang2025soft} was proposed with the same intuition.
It needs no training: at each thinking step it feeds the probability-weighted average of the top candidates'
embeddings in place of one sampled token.

Whether off-the-shelf models use these blends is contested.
Wu et al.~\cite{wu2025single} find that vanilla \ST{} is dominated by its top token and drifts into greedy, repetitive
decoding, which they call the Greedy Pitfall.
They add Dirichlet or Gumbel noise to the blend weights and report small accuracy gains over matched sampling.
Rizvi-Martel et al.~\cite{rizvimartel2026illusion} find no sign that vanilla \ST{} reasons in superposition,
and Backour~\cite{backour2026dynamics} shows that pretrained models collapse injected mixtures of embeddings through depth.

Gumbel noise also changes what the superposition question asks.
By the Gumbel-max trick \cite{gumbel1954statistical,maddison2014sampling}, the top token of a Gumbel blend is
distributed exactly as a sampled token.
\GST{} therefore feeds the token that sampling would feed, plus a \emph{residual}: the blend weight that sits on the
other candidates (\cref{fig:concept}).
If the model ignores the residual, \GST{} is sampling in disguise.
If the model carries both candidates forward and uses them, it reasons in genuine superposition.
A third outcome lies between the two: the model may read the residual without its answers depending on it.
Our research question is which of these holds:
\begin{quote}
\emph{Does \GST{} reason in superposition, or is it sampling in disguise?}
\end{quote}
We answer it at the level of single tokens.
Every blend in this paper, whether an exact 50/50 blend or one fed by \GST{}, is a \emph{forced} superposition,
built at the input, as opposed to a \emph{learned} one that emerges from training \cite{rizvimartel2026illusion}.
We ask whether one token position carries two candidates and whether that changes what the model does.
Whether whole reasoning paths are carried in superposition over many steps, as in the graph-search construction of
Zhu et al.~\cite{zhu2025superposition}, is a different question that we do not test.

We split the question into four measurements (\cref{fig:questions}) on DeepSeek-R1-Distill-Qwen-1.5B % SETUP
\cite{deepseekai2025r1,guo2025nature} solving GSM8K problems \cite{cobbe2021gsm8k}.
Q1 asks whether a blend of two candidates survives the layers.
Q2 asks whether information about the losing candidate, the runner-up, persists in the key-value cache and shifts
what later text predicts.
Q3 asks what blends the noise actually feeds in real runs.
Q4 asks whether the residual changes the model's next-token predictions and its final answers.
We committed numeric predictions for the main tests before running them.

Our evidence supports three statements about token-level superposition in this model.
First, the top token of \GST{} is a sampled token.
Second, the residual is read: it shifts the model's next-token predictions far more than a random residual of the
same size.
Third, no change in accuracy is detected on 100 problems, with a paired interval that still allows a few points % C23.1
either way.
The superposition is present in the input and read, but we find no evidence that it is used.
Our contributions are as follows.
\begin{itemize}[leftmargin=1.2em,itemsep=2pt,topsep=3pt]
\item \textbf{\GST{} is sampling plus a residual.}
We make the Gumbel-max consequence explicit and confirm it in logged runs:
at forks the top token changes at a rate of 0.45, the same as under sampling (0.45). % C22.4
Real forks feed lopsided blends, near one token at 47\% of forks, with a median residual of 0.12. % C22.11 C22.12
Dirichlet noise does not even match sampling: its top token leans toward the model's first choice
(\cref{sec:q3}).
\item \textbf{The runner-up persists in the cache.}
Blends collapse gradually through the layers, as Backour~\cite{backour2026dynamics} found in other models.
Yet information about the runner-up persists, mostly in the keys and values of layers 0 to 7, which close a pooled % C21.10
share of 0.71 to 0.91 of the readback gap (pre-registered); part of this is an echo of the input in the first layers. % C21.10
Later text that turns toward the runner-up shifts the projected next-token readout toward it ($+0.44$ for number % C21.7-8
vs number forks).
On the model's own path the shift is not detectable (\cref{sec:q1,sec:q2}).
\item \textbf{The model reads the residual, but accuracy does not change detectably.}
A coupled ablation runs sampling, \ST{} and a \ctrl{} in lockstep with shared noise.
\ST{} leaves sampling first on 59 problems against 24 for the control ($p = 0.0002$) % C23.5-6
and moves next-token predictions a median 3.9 times as much (95\% CI [3.1, 5.0]). % C23.15
Accuracy does not change detectably (88 vs 86 of 100; paired difference $+2$ points, 95\% CI [$-5$, $+9$]), % C23.1 C23.17
and paired by problem, the real and random residuals change sampling's final answer on mostly the same problems
(\cref{sec:q4}).
\item \textbf{A checkable study.}
Predictions were committed before the runs and the failed ones are reported as written (\cref{app:prereg}).
Every number from our experiments is recomputed from saved outputs by scripts in our repository (\cref{app:repro}).
\end{itemize}

\begin{figure*}[t]
\centering
% Fig 1: concept diagram (TikZ). In the paper:
%   \begin{figure*}[t] \centering % Fig 1: concept diagram (TikZ). In the paper:
%   \begin{figure*}[t] \centering % Fig 1: concept diagram (TikZ). In the paper:
%   \begin{figure*}[t] \centering \input{figures/fig1.tex} \caption{...} \label{fig:concept} \end{figure*}
% Preamble needs \usepackage{tikz} and \usetikzlibrary{arrows.meta,decorations.pathreplacing}.
% Drawn at final size (under the 6.5 in text width); include it as is, never scaled, so text stays 7-8 pt.
% Draft caption:
%   Gumbel Soft Thinking feeds the sampled token plus a residual. At each thinking step the model gives next-token
%   probabilities. Sampling feeds one token. Gumbel Soft Thinking adds Gumbel noise and feeds the softmax-weighted
%   mix of token embeddings; by the Gumbel-max trick, its top token is distributed exactly as the sampled one (and
%   is the same token when both use the same noise, as in Exp 23), so all it adds is the residual (blue). The
%   random-residual control keeps the top token and the residual's size but moves the residual to a random token
%   (orange). Weights are illustrative.
% No data numbers. Colours match figures/style.py: sampling grey, Gumbel Soft Thinking blue, control orange.
\definecolor{fsamp}{HTML}{7F7F7F}%
\definecolor{fsoft}{HTML}{0072B2}%
\definecolor{frand}{HTML}{E69F00}%
\definecolor{fdark}{HTML}{333333}%
\begin{tikzpicture}[x=1cm, y=1cm, font=\footnotesize, line width=0.5pt,
    >={Stealth[length=4.5pt, width=4pt]},
    wire/.style={line width=0.8pt, fdark},
    lab/.style={font=\scriptsize, inner sep=1pt, text=fdark},
    inbar/.style={font=\scriptsize\ttfamily, text=white, inner sep=0pt},
    brace/.style={decorate, decoration={brace, amplitude=3.5pt}, black!55, line width=0.5pt}]
\def\L{5.5}            % bar length = weight 1
\def\xb{7.45}           % bar start
\pgfmathsetmacro{\xt}{\xb + 0.75*\L}   % end of the top token (weight 0.75)
\pgfmathsetmacro{\xe}{\xb + \L}

% ---- next-token probabilities ----
\node[lab, align=center, anchor=south] at (1.45, 2.95) {next-token\\[-1pt]probabilities};
\foreach \x/\p/\t in {0.8/0.55/7, 1.45/0.35/8, 2.1/0.10/12} {
  \draw[fdark, fill=white] (\x-0.21, 0.6) rectangle (\x+0.21, 0.6+3.9*\p*0.95);
  \node[lab, font=\scriptsize\ttfamily, anchor=north] at (\x, 0.55) {\t};
}
\draw[fdark] (0.4, 0.6) -- (2.5, 0.6);

% ---- three ways to feed it ----
\draw[wire] (2.65, 1.975) -- (3.2, 1.975);
\draw[wire] (3.2, 0.725) -- (3.2, 3.225);
\foreach \y/\t in {3.225/{sample one token}, 1.975/{add Gumbel noise, softmax}, 0.725/{same, then move the residual}} {
  \draw[wire, ->] (3.2, \y) -- (\xb-0.08, \y);
  \node[lab, anchor=south] at ({(3.2+\xb)/2}, \y+0.03) {\t};
}
\fill[fdark] (3.2, 1.975) circle (1.3pt);

% ---- what is fed: one bar per method, length = total weight 1 ----
\fill[fsamp] (\xb, 3.0) rectangle (\xe, 3.45);
\node[inbar] at ({(\xb+\xt)/2}, 3.225) {7};
\fill[fsamp] (\xb, 1.75) rectangle (\xt, 2.2);
\fill[fsoft] (\xt, 1.75) rectangle ({\xb+0.95*\L}, 2.2);
\fill[fsoft!55] ({\xb+0.95*\L}, 1.75) rectangle (\xe, 2.2);
\draw[white, line width=0.8pt] ({\xb+0.95*\L}, 1.75) -- ({\xb+0.95*\L}, 2.2);
\node[inbar] at ({(\xb+\xt)/2}, 1.975) {7};
\node[inbar] at ({\xt+0.1*\L}, 1.975) {8};
\fill[fsamp] (\xb, 0.5) rectangle (\xt, 0.95);
\fill[frand] (\xt, 0.5) rectangle (\xe, 0.95);
\node[inbar] at ({(\xb+\xt)/2}, 0.725) {7};
\draw[black!35, dash pattern=on 2pt off 1.5pt] (\xt, 0.35) -- (\xt, 3.6);

% ---- braces: the shared top token and the residual ----
\draw[brace] (\xb, 3.62) -- (\xt-0.04, 3.62) node[midway, above=4pt, lab] {top token (the same in all three)};
\draw[brace] (\xt+0.04, 3.62) -- (\xe, 3.62) node[midway, above=4pt, lab] {residual};

% ---- names ----
\foreach \y/\c/\n/\d in {3.225/fsamp/{sampling}/{one token}, 1.975/fsoft/{Gumbel Soft Thinking}/{token $+$ residual},
                         0.725/frand/{random-residual control}/{residual on a random token}} {
  \node[lab, anchor=west, align=left, text=\c] at (\xe+0.18, \y) {\textbf{\n}\\[-1pt]\textcolor{fdark}{\d}};
}
\end{tikzpicture}
 \caption{...} \label{fig:concept} \end{figure*}
% Preamble needs \usepackage{tikz} and \usetikzlibrary{arrows.meta,decorations.pathreplacing}.
% Drawn at final size (under the 6.5 in text width); include it as is, never scaled, so text stays 7-8 pt.
% Draft caption:
%   Gumbel Soft Thinking feeds the sampled token plus a residual. At each thinking step the model gives next-token
%   probabilities. Sampling feeds one token. Gumbel Soft Thinking adds Gumbel noise and feeds the softmax-weighted
%   mix of token embeddings; by the Gumbel-max trick, its top token is distributed exactly as the sampled one (and
%   is the same token when both use the same noise, as in Exp 23), so all it adds is the residual (blue). The
%   random-residual control keeps the top token and the residual's size but moves the residual to a random token
%   (orange). Weights are illustrative.
% No data numbers. Colours match figures/style.py: sampling grey, Gumbel Soft Thinking blue, control orange.
\definecolor{fsamp}{HTML}{7F7F7F}%
\definecolor{fsoft}{HTML}{0072B2}%
\definecolor{frand}{HTML}{E69F00}%
\definecolor{fdark}{HTML}{333333}%
\begin{tikzpicture}[x=1cm, y=1cm, font=\footnotesize, line width=0.5pt,
    >={Stealth[length=4.5pt, width=4pt]},
    wire/.style={line width=0.8pt, fdark},
    lab/.style={font=\scriptsize, inner sep=1pt, text=fdark},
    inbar/.style={font=\scriptsize\ttfamily, text=white, inner sep=0pt},
    brace/.style={decorate, decoration={brace, amplitude=3.5pt}, black!55, line width=0.5pt}]
\def\L{5.5}            % bar length = weight 1
\def\xb{7.45}           % bar start
\pgfmathsetmacro{\xt}{\xb + 0.75*\L}   % end of the top token (weight 0.75)
\pgfmathsetmacro{\xe}{\xb + \L}

% ---- next-token probabilities ----
\node[lab, align=center, anchor=south] at (1.45, 2.95) {next-token\\[-1pt]probabilities};
\foreach \x/\p/\t in {0.8/0.55/7, 1.45/0.35/8, 2.1/0.10/12} {
  \draw[fdark, fill=white] (\x-0.21, 0.6) rectangle (\x+0.21, 0.6+3.9*\p*0.95);
  \node[lab, font=\scriptsize\ttfamily, anchor=north] at (\x, 0.55) {\t};
}
\draw[fdark] (0.4, 0.6) -- (2.5, 0.6);

% ---- three ways to feed it ----
\draw[wire] (2.65, 1.975) -- (3.2, 1.975);
\draw[wire] (3.2, 0.725) -- (3.2, 3.225);
\foreach \y/\t in {3.225/{sample one token}, 1.975/{add Gumbel noise, softmax}, 0.725/{same, then move the residual}} {
  \draw[wire, ->] (3.2, \y) -- (\xb-0.08, \y);
  \node[lab, anchor=south] at ({(3.2+\xb)/2}, \y+0.03) {\t};
}
\fill[fdark] (3.2, 1.975) circle (1.3pt);

% ---- what is fed: one bar per method, length = total weight 1 ----
\fill[fsamp] (\xb, 3.0) rectangle (\xe, 3.45);
\node[inbar] at ({(\xb+\xt)/2}, 3.225) {7};
\fill[fsamp] (\xb, 1.75) rectangle (\xt, 2.2);
\fill[fsoft] (\xt, 1.75) rectangle ({\xb+0.95*\L}, 2.2);
\fill[fsoft!55] ({\xb+0.95*\L}, 1.75) rectangle (\xe, 2.2);
\draw[white, line width=0.8pt] ({\xb+0.95*\L}, 1.75) -- ({\xb+0.95*\L}, 2.2);
\node[inbar] at ({(\xb+\xt)/2}, 1.975) {7};
\node[inbar] at ({\xt+0.1*\L}, 1.975) {8};
\fill[fsamp] (\xb, 0.5) rectangle (\xt, 0.95);
\fill[frand] (\xt, 0.5) rectangle (\xe, 0.95);
\node[inbar] at ({(\xb+\xt)/2}, 0.725) {7};
\draw[black!35, dash pattern=on 2pt off 1.5pt] (\xt, 0.35) -- (\xt, 3.6);

% ---- braces: the shared top token and the residual ----
\draw[brace] (\xb, 3.62) -- (\xt-0.04, 3.62) node[midway, above=4pt, lab] {top token (the same in all three)};
\draw[brace] (\xt+0.04, 3.62) -- (\xe, 3.62) node[midway, above=4pt, lab] {residual};

% ---- names ----
\foreach \y/\c/\n/\d in {3.225/fsamp/{sampling}/{one token}, 1.975/fsoft/{Gumbel Soft Thinking}/{token $+$ residual},
                         0.725/frand/{random-residual control}/{residual on a random token}} {
  \node[lab, anchor=west, align=left, text=\c] at (\xe+0.18, \y) {\textbf{\n}\\[-1pt]\textcolor{fdark}{\d}};
}
\end{tikzpicture}
 \caption{...} \label{fig:concept} \end{figure*}
% Preamble needs \usepackage{tikz} and \usetikzlibrary{arrows.meta,decorations.pathreplacing}.
% Drawn at final size (under the 6.5 in text width); include it as is, never scaled, so text stays 7-8 pt.
% Draft caption:
%   Gumbel Soft Thinking feeds the sampled token plus a residual. At each thinking step the model gives next-token
%   probabilities. Sampling feeds one token. Gumbel Soft Thinking adds Gumbel noise and feeds the softmax-weighted
%   mix of token embeddings; by the Gumbel-max trick, its top token is distributed exactly as the sampled one (and
%   is the same token when both use the same noise, as in Exp 23), so all it adds is the residual (blue). The
%   random-residual control keeps the top token and the residual's size but moves the residual to a random token
%   (orange). Weights are illustrative.
% No data numbers. Colours match figures/style.py: sampling grey, Gumbel Soft Thinking blue, control orange.
\definecolor{fsamp}{HTML}{7F7F7F}%
\definecolor{fsoft}{HTML}{0072B2}%
\definecolor{frand}{HTML}{E69F00}%
\definecolor{fdark}{HTML}{333333}%
\begin{tikzpicture}[x=1cm, y=1cm, font=\footnotesize, line width=0.5pt,
    >={Stealth[length=4.5pt, width=4pt]},
    wire/.style={line width=0.8pt, fdark},
    lab/.style={font=\scriptsize, inner sep=1pt, text=fdark},
    inbar/.style={font=\scriptsize\ttfamily, text=white, inner sep=0pt},
    brace/.style={decorate, decoration={brace, amplitude=3.5pt}, black!55, line width=0.5pt}]
\def\L{5.5}            % bar length = weight 1
\def\xb{7.45}           % bar start
\pgfmathsetmacro{\xt}{\xb + 0.75*\L}   % end of the top token (weight 0.75)
\pgfmathsetmacro{\xe}{\xb + \L}

% ---- next-token probabilities ----
\node[lab, align=center, anchor=south] at (1.45, 2.95) {next-token\\[-1pt]probabilities};
\foreach \x/\p/\t in {0.8/0.55/7, 1.45/0.35/8, 2.1/0.10/12} {
  \draw[fdark, fill=white] (\x-0.21, 0.6) rectangle (\x+0.21, 0.6+3.9*\p*0.95);
  \node[lab, font=\scriptsize\ttfamily, anchor=north] at (\x, 0.55) {\t};
}
\draw[fdark] (0.4, 0.6) -- (2.5, 0.6);

% ---- three ways to feed it ----
\draw[wire] (2.65, 1.975) -- (3.2, 1.975);
\draw[wire] (3.2, 0.725) -- (3.2, 3.225);
\foreach \y/\t in {3.225/{sample one token}, 1.975/{add Gumbel noise, softmax}, 0.725/{same, then move the residual}} {
  \draw[wire, ->] (3.2, \y) -- (\xb-0.08, \y);
  \node[lab, anchor=south] at ({(3.2+\xb)/2}, \y+0.03) {\t};
}
\fill[fdark] (3.2, 1.975) circle (1.3pt);

% ---- what is fed: one bar per method, length = total weight 1 ----
\fill[fsamp] (\xb, 3.0) rectangle (\xe, 3.45);
\node[inbar] at ({(\xb+\xt)/2}, 3.225) {7};
\fill[fsamp] (\xb, 1.75) rectangle (\xt, 2.2);
\fill[fsoft] (\xt, 1.75) rectangle ({\xb+0.95*\L}, 2.2);
\fill[fsoft!55] ({\xb+0.95*\L}, 1.75) rectangle (\xe, 2.2);
\draw[white, line width=0.8pt] ({\xb+0.95*\L}, 1.75) -- ({\xb+0.95*\L}, 2.2);
\node[inbar] at ({(\xb+\xt)/2}, 1.975) {7};
\node[inbar] at ({\xt+0.1*\L}, 1.975) {8};
\fill[fsamp] (\xb, 0.5) rectangle (\xt, 0.95);
\fill[frand] (\xt, 0.5) rectangle (\xe, 0.95);
\node[inbar] at ({(\xb+\xt)/2}, 0.725) {7};
\draw[black!35, dash pattern=on 2pt off 1.5pt] (\xt, 0.35) -- (\xt, 3.6);

% ---- braces: the shared top token and the residual ----
\draw[brace] (\xb, 3.62) -- (\xt-0.04, 3.62) node[midway, above=4pt, lab] {top token (the same in all three)};
\draw[brace] (\xt+0.04, 3.62) -- (\xe, 3.62) node[midway, above=4pt, lab] {residual};

% ---- names ----
\foreach \y/\c/\n/\d in {3.225/fsamp/{sampling}/{one token}, 1.975/fsoft/{Gumbel Soft Thinking}/{token $+$ residual},
                         0.725/frand/{random-residual control}/{residual on a random token}} {
  \node[lab, anchor=west, align=left, text=\c] at (\xe+0.18, \y) {\textbf{\n}\\[-1pt]\textcolor{fdark}{\d}};
}
\end{tikzpicture}

\caption{\GST{} feeds the sampled token plus a residual.
At each thinking step the model gives next-token probabilities.
Sampling feeds one token.
\GST{} adds Gumbel noise and feeds the softmax-weighted mix of token embeddings.
By the Gumbel-max trick, its top token is distributed exactly as the sampled one,
and it is the same token when both use the same noise in the same context, as in our coupled ablation (\cref{sec:q4}).
So all it adds is the residual (blue).
The \ctrl{} keeps the top token and the residual's size but moves the residual to a random token (orange).
Weights are illustrative.}
\label{fig:concept}
\end{figure*}

\begin{figure*}[t]
\centering
% Fig 2: the four questions (TikZ). Follows Fig 1 (the concept). In the paper:
%   \begin{figure*}[t] \centering % Fig 2: the four questions (TikZ). Follows Fig 1 (the concept). In the paper:
%   \begin{figure*}[t] \centering % Fig 2: the four questions (TikZ). Follows Fig 1 (the concept). In the paper:
%   \begin{figure*}[t] \centering \input{figures/fig2.tex} \caption{...} \label{fig:questions} \end{figure*}
% Preamble needs \usepackage{tikz} and \usetikzlibrary{arrows.meta}. Drawn at final size; include it as is.
% Draft caption:
%   The four questions we test. (1) Does a blend of two candidates survive the layers, or snap to one of them
%   (Fig. 3)? (2) Is the runner-up stored in the K,V cache, and can later steps read it back (Figs. 4-5)?
%   (3) What blends does the noise actually feed in real runs (Fig. 6)? (4) Does the residual move the model's
%   next-token predictions more than a random residual of the same size (Fig. 7, Table 1)? Drawings are schematic.
% No data numbers: every position and bar is illustrative.
% Colours match figures/style.py: sampling grey, Gumbel Soft Thinking blue, random-residual control orange.
\definecolor{fsamp}{HTML}{7F7F7F}%
\definecolor{fsoft}{HTML}{0072B2}%
\definecolor{frand}{HTML}{E69F00}%
\definecolor{fdark}{HTML}{333333}%
\definecolor{fcell}{HTML}{E9E9EC}%
\begin{tikzpicture}[x=1cm, y=1cm, font=\footnotesize, line width=0.5pt,
    >={Stealth[length=4pt, width=3.5pt]},
    lab/.style={font=\scriptsize, inner sep=1pt, text=black!60},
    qm/.style={circle, fill=fdark, text=white, font=\scriptsize\bfseries, inner sep=0pt, minimum size=10.5pt},
    qt/.style={anchor=north west, font=\footnotesize\bfseries, text=fdark, inner sep=0pt, text width=3.4cm,
               align=left, execute at begin node={\hyphenpenalty=10000}}]

\foreach \x in {4.075, 8.225, 12.375} \draw[black!15] (\x, 0.05) -- (\x, 3.65);

% ---------- 1: does the blend survive the layers? ----------
\begin{scope}[shift={(0, 0)}]
  \node[qm] at (0.17, 3.5) {1}; \node[qt] at (0.42, 3.6) {Does the blend survive the layers?};
  \foreach \y/\f in {0.6/0.5, 1.45/0.72, 2.3/0.87} {
    \draw[black!30, line width=0.8pt] (1.0, \y) -- (3.6, \y);
    \draw[fdark, fill=white] (1.0, \y) circle (0.075);
    \fill[fsamp] (3.6, \y) circle (0.075);
    \fill[fsoft] (1.0+2.6*\f, \y) circle (0.095);
  }
  \draw[->, fsoft!60, dash pattern=on 1.5pt off 1.2pt] (2.38, 0.72) .. controls (2.52, 1.05) .. (2.79, 1.34);
  \draw[->, fsoft!60, dash pattern=on 1.5pt off 1.2pt] (2.95, 1.57) .. controls (3.08, 1.9) .. (3.22, 2.18);
  \node[lab, anchor=east] at (0.85, 0.6) {input};
  \node[lab, anchor=east] at (0.85, 2.3) {output};
  \node[lab, font=\scriptsize\ttfamily, text=fdark] at (1.0, 0.3) {8};
  \node[lab, font=\scriptsize\ttfamily, text=fdark] at (3.6, 0.3) {7};
  \node[lab, text=fsoft] at (2.3, 0.3) {blend};
\end{scope}

% ---------- 2: is the runner-up stored? ----------
\begin{scope}[shift={(4.15, 0)}]
  \node[qm] at (0.17, 3.5) {2}; \node[qt] at (0.42, 3.6) {Does the runner-up persist?};
  \foreach \i in {0,...,4} \foreach \j in {0,...,3}
    \fill[fcell] (0.75+0.57*\i, 0.5+0.45*\j) rectangle (1.25+0.57*\i, 0.88+0.45*\j);
  \fill[fsoft] (1.32, 0.5) rectangle (1.82, 0.88);
  \fill[fsoft!60] (1.32, 0.95) rectangle (1.82, 1.33);
  \fill[fsoft!22] (1.32, 1.4) rectangle (1.82, 1.78);
  \draw[->, fdark, line width=0.7pt] (3.28, 2.3) .. controls (3.15, 2.85) and (1.75, 2.85) .. (1.6, 2.32);
  \node[lab, text=fdark] at (2.44, 2.92) {read later};
  \draw[->, black!45] (0.5, 0.5) -- (0.5, 2.23) node[midway, lab, rotate=90, anchor=south] {layers};
  \node[lab, text=fsoft] at (1.57, 0.27) {fork};
  \draw[->, black!45] (2.1, 0.27) -- (3.0, 0.27);
  \node[lab, anchor=west] at (3.03, 0.27) {steps};
\end{scope}

% ---------- 3: what does the noise feed? ----------
\begin{scope}[shift={(8.3, 0)}]
  \node[qm] at (0.17, 3.5) {3}; \node[qt] at (0.42, 3.6) {What does the noise feed?};
  \foreach \r [count=\k] in {0.04, 0.15, 0.02, 0.38, 0.07, 0.01, 0.26, 0.05} {
    \pgfmathsetmacro{\yy}{2.55-0.27*\k}
    \fill[fsamp] (0.55, \yy) rectangle ({0.55+3.0*(1-\r)}, \yy+0.17);
    \fill[fsoft] ({0.55+3.0*(1-\r)}, \yy) rectangle (3.55, \yy+0.17);
  }
  \node[lab] at (2.05, 0.12) {fed blends in real runs};
\end{scope}

% ---------- 4: is the residual used? ----------
\begin{scope}[shift={(12.45, 0)}]
  \node[qm] at (0.17, 3.5) {4}; \node[qt] at (0.42, 3.6) {Is the residual used?};
  \draw[fill=fcell, draw=black!30, rounded corners=2pt] (1.4, 0.3) rectangle (1.95, 2.35);
  \node[lab, rotate=90, text=fdark] at (1.675, 1.325) {model};
  \foreach \y/\c/\h in {2.05/fsamp/{0.48,0.30,0.14,0.06}, 1.3/fsoft/{0.30,0.46,0.14,0.08}, 0.55/frand/{0.46,0.31,0.15,0.07}} {
    \fill[fsamp] (0.1, \y-0.14) rectangle (0.74, \y+0.14);
    \draw[->, black!45] (1.03, \y) -- (1.34, \y);
    \draw[->, black!45] (2.0, \y) -- (2.3, \y);
    \foreach \v [count=\k] in \h \fill[\c] ({2.38+0.22*(\k-1)}, \y-0.22) rectangle ({2.53+0.22*(\k-1)}, {\y-0.22+0.85*\v});
    \draw[black!40] (2.33, \y-0.22) -- (3.25, \y-0.22);
  }
  \fill[fsamp] (0.74, 1.91) rectangle (0.95, 2.19);
  \fill[fsoft] (0.74, 1.16) rectangle (0.95, 1.44);
  \fill[frand] (0.74, 0.41) rectangle (0.95, 0.69);
  \foreach \y in {1.3, 0.55} \foreach \v [count=\k] in {0.48,0.30,0.14,0.06}
    \draw[fdark, dash pattern=on 1pt off 0.8pt, line width=0.35pt] ({2.38+0.22*(\k-1)}, \y-0.22) rectangle ({2.53+0.22*(\k-1)}, {\y-0.22+0.85*\v});
  \node[lab] at (2.79, 2.55) {next token};
\end{scope}
\end{tikzpicture}
 \caption{...} \label{fig:questions} \end{figure*}
% Preamble needs \usepackage{tikz} and \usetikzlibrary{arrows.meta}. Drawn at final size; include it as is.
% Draft caption:
%   The four questions we test. (1) Does a blend of two candidates survive the layers, or snap to one of them
%   (Fig. 3)? (2) Is the runner-up stored in the K,V cache, and can later steps read it back (Figs. 4-5)?
%   (3) What blends does the noise actually feed in real runs (Fig. 6)? (4) Does the residual move the model's
%   next-token predictions more than a random residual of the same size (Fig. 7, Table 1)? Drawings are schematic.
% No data numbers: every position and bar is illustrative.
% Colours match figures/style.py: sampling grey, Gumbel Soft Thinking blue, random-residual control orange.
\definecolor{fsamp}{HTML}{7F7F7F}%
\definecolor{fsoft}{HTML}{0072B2}%
\definecolor{frand}{HTML}{E69F00}%
\definecolor{fdark}{HTML}{333333}%
\definecolor{fcell}{HTML}{E9E9EC}%
\begin{tikzpicture}[x=1cm, y=1cm, font=\footnotesize, line width=0.5pt,
    >={Stealth[length=4pt, width=3.5pt]},
    lab/.style={font=\scriptsize, inner sep=1pt, text=black!60},
    qm/.style={circle, fill=fdark, text=white, font=\scriptsize\bfseries, inner sep=0pt, minimum size=10.5pt},
    qt/.style={anchor=north west, font=\footnotesize\bfseries, text=fdark, inner sep=0pt, text width=3.4cm,
               align=left, execute at begin node={\hyphenpenalty=10000}}]

\foreach \x in {4.075, 8.225, 12.375} \draw[black!15] (\x, 0.05) -- (\x, 3.65);

% ---------- 1: does the blend survive the layers? ----------
\begin{scope}[shift={(0, 0)}]
  \node[qm] at (0.17, 3.5) {1}; \node[qt] at (0.42, 3.6) {Does the blend survive the layers?};
  \foreach \y/\f in {0.6/0.5, 1.45/0.72, 2.3/0.87} {
    \draw[black!30, line width=0.8pt] (1.0, \y) -- (3.6, \y);
    \draw[fdark, fill=white] (1.0, \y) circle (0.075);
    \fill[fsamp] (3.6, \y) circle (0.075);
    \fill[fsoft] (1.0+2.6*\f, \y) circle (0.095);
  }
  \draw[->, fsoft!60, dash pattern=on 1.5pt off 1.2pt] (2.38, 0.72) .. controls (2.52, 1.05) .. (2.79, 1.34);
  \draw[->, fsoft!60, dash pattern=on 1.5pt off 1.2pt] (2.95, 1.57) .. controls (3.08, 1.9) .. (3.22, 2.18);
  \node[lab, anchor=east] at (0.85, 0.6) {input};
  \node[lab, anchor=east] at (0.85, 2.3) {output};
  \node[lab, font=\scriptsize\ttfamily, text=fdark] at (1.0, 0.3) {8};
  \node[lab, font=\scriptsize\ttfamily, text=fdark] at (3.6, 0.3) {7};
  \node[lab, text=fsoft] at (2.3, 0.3) {blend};
\end{scope}

% ---------- 2: is the runner-up stored? ----------
\begin{scope}[shift={(4.15, 0)}]
  \node[qm] at (0.17, 3.5) {2}; \node[qt] at (0.42, 3.6) {Does the runner-up persist?};
  \foreach \i in {0,...,4} \foreach \j in {0,...,3}
    \fill[fcell] (0.75+0.57*\i, 0.5+0.45*\j) rectangle (1.25+0.57*\i, 0.88+0.45*\j);
  \fill[fsoft] (1.32, 0.5) rectangle (1.82, 0.88);
  \fill[fsoft!60] (1.32, 0.95) rectangle (1.82, 1.33);
  \fill[fsoft!22] (1.32, 1.4) rectangle (1.82, 1.78);
  \draw[->, fdark, line width=0.7pt] (3.28, 2.3) .. controls (3.15, 2.85) and (1.75, 2.85) .. (1.6, 2.32);
  \node[lab, text=fdark] at (2.44, 2.92) {read later};
  \draw[->, black!45] (0.5, 0.5) -- (0.5, 2.23) node[midway, lab, rotate=90, anchor=south] {layers};
  \node[lab, text=fsoft] at (1.57, 0.27) {fork};
  \draw[->, black!45] (2.1, 0.27) -- (3.0, 0.27);
  \node[lab, anchor=west] at (3.03, 0.27) {steps};
\end{scope}

% ---------- 3: what does the noise feed? ----------
\begin{scope}[shift={(8.3, 0)}]
  \node[qm] at (0.17, 3.5) {3}; \node[qt] at (0.42, 3.6) {What does the noise feed?};
  \foreach \r [count=\k] in {0.04, 0.15, 0.02, 0.38, 0.07, 0.01, 0.26, 0.05} {
    \pgfmathsetmacro{\yy}{2.55-0.27*\k}
    \fill[fsamp] (0.55, \yy) rectangle ({0.55+3.0*(1-\r)}, \yy+0.17);
    \fill[fsoft] ({0.55+3.0*(1-\r)}, \yy) rectangle (3.55, \yy+0.17);
  }
  \node[lab] at (2.05, 0.12) {fed blends in real runs};
\end{scope}

% ---------- 4: is the residual used? ----------
\begin{scope}[shift={(12.45, 0)}]
  \node[qm] at (0.17, 3.5) {4}; \node[qt] at (0.42, 3.6) {Is the residual used?};
  \draw[fill=fcell, draw=black!30, rounded corners=2pt] (1.4, 0.3) rectangle (1.95, 2.35);
  \node[lab, rotate=90, text=fdark] at (1.675, 1.325) {model};
  \foreach \y/\c/\h in {2.05/fsamp/{0.48,0.30,0.14,0.06}, 1.3/fsoft/{0.30,0.46,0.14,0.08}, 0.55/frand/{0.46,0.31,0.15,0.07}} {
    \fill[fsamp] (0.1, \y-0.14) rectangle (0.74, \y+0.14);
    \draw[->, black!45] (1.03, \y) -- (1.34, \y);
    \draw[->, black!45] (2.0, \y) -- (2.3, \y);
    \foreach \v [count=\k] in \h \fill[\c] ({2.38+0.22*(\k-1)}, \y-0.22) rectangle ({2.53+0.22*(\k-1)}, {\y-0.22+0.85*\v});
    \draw[black!40] (2.33, \y-0.22) -- (3.25, \y-0.22);
  }
  \fill[fsamp] (0.74, 1.91) rectangle (0.95, 2.19);
  \fill[fsoft] (0.74, 1.16) rectangle (0.95, 1.44);
  \fill[frand] (0.74, 0.41) rectangle (0.95, 0.69);
  \foreach \y in {1.3, 0.55} \foreach \v [count=\k] in {0.48,0.30,0.14,0.06}
    \draw[fdark, dash pattern=on 1pt off 0.8pt, line width=0.35pt] ({2.38+0.22*(\k-1)}, \y-0.22) rectangle ({2.53+0.22*(\k-1)}, {\y-0.22+0.85*\v});
  \node[lab] at (2.79, 2.55) {next token};
\end{scope}
\end{tikzpicture}
 \caption{...} \label{fig:questions} \end{figure*}
% Preamble needs \usepackage{tikz} and \usetikzlibrary{arrows.meta}. Drawn at final size; include it as is.
% Draft caption:
%   The four questions we test. (1) Does a blend of two candidates survive the layers, or snap to one of them
%   (Fig. 3)? (2) Is the runner-up stored in the K,V cache, and can later steps read it back (Figs. 4-5)?
%   (3) What blends does the noise actually feed in real runs (Fig. 6)? (4) Does the residual move the model's
%   next-token predictions more than a random residual of the same size (Fig. 7, Table 1)? Drawings are schematic.
% No data numbers: every position and bar is illustrative.
% Colours match figures/style.py: sampling grey, Gumbel Soft Thinking blue, random-residual control orange.
\definecolor{fsamp}{HTML}{7F7F7F}%
\definecolor{fsoft}{HTML}{0072B2}%
\definecolor{frand}{HTML}{E69F00}%
\definecolor{fdark}{HTML}{333333}%
\definecolor{fcell}{HTML}{E9E9EC}%
\begin{tikzpicture}[x=1cm, y=1cm, font=\footnotesize, line width=0.5pt,
    >={Stealth[length=4pt, width=3.5pt]},
    lab/.style={font=\scriptsize, inner sep=1pt, text=black!60},
    qm/.style={circle, fill=fdark, text=white, font=\scriptsize\bfseries, inner sep=0pt, minimum size=10.5pt},
    qt/.style={anchor=north west, font=\footnotesize\bfseries, text=fdark, inner sep=0pt, text width=3.4cm,
               align=left, execute at begin node={\hyphenpenalty=10000}}]

\foreach \x in {4.075, 8.225, 12.375} \draw[black!15] (\x, 0.05) -- (\x, 3.65);

% ---------- 1: does the blend survive the layers? ----------
\begin{scope}[shift={(0, 0)}]
  \node[qm] at (0.17, 3.5) {1}; \node[qt] at (0.42, 3.6) {Does the blend survive the layers?};
  \foreach \y/\f in {0.6/0.5, 1.45/0.72, 2.3/0.87} {
    \draw[black!30, line width=0.8pt] (1.0, \y) -- (3.6, \y);
    \draw[fdark, fill=white] (1.0, \y) circle (0.075);
    \fill[fsamp] (3.6, \y) circle (0.075);
    \fill[fsoft] (1.0+2.6*\f, \y) circle (0.095);
  }
  \draw[->, fsoft!60, dash pattern=on 1.5pt off 1.2pt] (2.38, 0.72) .. controls (2.52, 1.05) .. (2.79, 1.34);
  \draw[->, fsoft!60, dash pattern=on 1.5pt off 1.2pt] (2.95, 1.57) .. controls (3.08, 1.9) .. (3.22, 2.18);
  \node[lab, anchor=east] at (0.85, 0.6) {input};
  \node[lab, anchor=east] at (0.85, 2.3) {output};
  \node[lab, font=\scriptsize\ttfamily, text=fdark] at (1.0, 0.3) {8};
  \node[lab, font=\scriptsize\ttfamily, text=fdark] at (3.6, 0.3) {7};
  \node[lab, text=fsoft] at (2.3, 0.3) {blend};
\end{scope}

% ---------- 2: is the runner-up stored? ----------
\begin{scope}[shift={(4.15, 0)}]
  \node[qm] at (0.17, 3.5) {2}; \node[qt] at (0.42, 3.6) {Does the runner-up persist?};
  \foreach \i in {0,...,4} \foreach \j in {0,...,3}
    \fill[fcell] (0.75+0.57*\i, 0.5+0.45*\j) rectangle (1.25+0.57*\i, 0.88+0.45*\j);
  \fill[fsoft] (1.32, 0.5) rectangle (1.82, 0.88);
  \fill[fsoft!60] (1.32, 0.95) rectangle (1.82, 1.33);
  \fill[fsoft!22] (1.32, 1.4) rectangle (1.82, 1.78);
  \draw[->, fdark, line width=0.7pt] (3.28, 2.3) .. controls (3.15, 2.85) and (1.75, 2.85) .. (1.6, 2.32);
  \node[lab, text=fdark] at (2.44, 2.92) {read later};
  \draw[->, black!45] (0.5, 0.5) -- (0.5, 2.23) node[midway, lab, rotate=90, anchor=south] {layers};
  \node[lab, text=fsoft] at (1.57, 0.27) {fork};
  \draw[->, black!45] (2.1, 0.27) -- (3.0, 0.27);
  \node[lab, anchor=west] at (3.03, 0.27) {steps};
\end{scope}

% ---------- 3: what does the noise feed? ----------
\begin{scope}[shift={(8.3, 0)}]
  \node[qm] at (0.17, 3.5) {3}; \node[qt] at (0.42, 3.6) {What does the noise feed?};
  \foreach \r [count=\k] in {0.04, 0.15, 0.02, 0.38, 0.07, 0.01, 0.26, 0.05} {
    \pgfmathsetmacro{\yy}{2.55-0.27*\k}
    \fill[fsamp] (0.55, \yy) rectangle ({0.55+3.0*(1-\r)}, \yy+0.17);
    \fill[fsoft] ({0.55+3.0*(1-\r)}, \yy) rectangle (3.55, \yy+0.17);
  }
  \node[lab] at (2.05, 0.12) {fed blends in real runs};
\end{scope}

% ---------- 4: is the residual used? ----------
\begin{scope}[shift={(12.45, 0)}]
  \node[qm] at (0.17, 3.5) {4}; \node[qt] at (0.42, 3.6) {Is the residual used?};
  \draw[fill=fcell, draw=black!30, rounded corners=2pt] (1.4, 0.3) rectangle (1.95, 2.35);
  \node[lab, rotate=90, text=fdark] at (1.675, 1.325) {model};
  \foreach \y/\c/\h in {2.05/fsamp/{0.48,0.30,0.14,0.06}, 1.3/fsoft/{0.30,0.46,0.14,0.08}, 0.55/frand/{0.46,0.31,0.15,0.07}} {
    \fill[fsamp] (0.1, \y-0.14) rectangle (0.74, \y+0.14);
    \draw[->, black!45] (1.03, \y) -- (1.34, \y);
    \draw[->, black!45] (2.0, \y) -- (2.3, \y);
    \foreach \v [count=\k] in \h \fill[\c] ({2.38+0.22*(\k-1)}, \y-0.22) rectangle ({2.53+0.22*(\k-1)}, {\y-0.22+0.85*\v});
    \draw[black!40] (2.33, \y-0.22) -- (3.25, \y-0.22);
  }
  \fill[fsamp] (0.74, 1.91) rectangle (0.95, 2.19);
  \fill[fsoft] (0.74, 1.16) rectangle (0.95, 1.44);
  \fill[frand] (0.74, 0.41) rectangle (0.95, 0.69);
  \foreach \y in {1.3, 0.55} \foreach \v [count=\k] in {0.48,0.30,0.14,0.06}
    \draw[fdark, dash pattern=on 1pt off 0.8pt, line width=0.35pt] ({2.38+0.22*(\k-1)}, \y-0.22) rectangle ({2.53+0.22*(\k-1)}, {\y-0.22+0.85*\v});
  \node[lab] at (2.79, 2.55) {next token};
\end{scope}
\end{tikzpicture}

\caption{The four questions we test.
(1) Does a blend of two candidates survive the layers, or snap to one of them (\cref{fig:snap})?
(2) Does information about the runner-up persist in the K,V cache and shift what later steps predict (\cref{fig:readback,fig:patch})?
(3) What blends does the noise actually feed in real runs (\cref{fig:fed})?
(4) Does the residual move the model's next-token predictions more than a random residual of the same size
(\cref{fig:ablation}, \cref{tab:accuracy})?
Drawings are schematic.}
\label{fig:questions}
\end{figure*}
